\documentclass[11pt]{report}

\begin{document}

\begin{center}
\textbf{Response to Peer Reviewers}
\end{center}

September 18, 2026

Dear reviewers,

Thank you to the reviewers for their careful reading of our manuscript, APOR-D-26-00768,
``Development, Validation, and Benchmarking of a Multidisciplinary Semi-Analytical Model for
Wave Energy Converters'' and for their constructive and thoughtful feedback.

We have considered all of the comments carefully and have revised the manuscript accordingly.
Our responses to each comment are provided below, with the corresponding changes to the manuscript described where applicable.

We believe these revisions have strengthened the manuscript and we look forward to your feedback.

Sincerely,

Rebecca McCabe, PhD

rgm222@cornell.edu

(On behalf of all authors)

\clearpage

\begin{enumerate}

% ----------------------------------------------------------------
% Reviewer 1
% ----------------------------------------------------------------

\item \textbf{Reviewer 1, Comment 1.}

\emph{The manuscript is excessively long and should be significantly condensed to improve readability.}

\textbf{Response:}
The manuscript has been substantially reorganized and condensed. See the response to Reviewer 1, Comment 2 below.

\item \textbf{Reviewer 1, Comment 2.}

\emph{The manuscript would benefit from improved organization, with redundant descriptions removed and clearer transitions between sections.}

\textbf{Response:}
The organization has been updated to focus on model structure, assumptions, verification, benchmarking, discussion, and limitations. Other content has been moved to the appendix.

\item \textbf{Reviewer 1, Comment 3.}

\emph{Validation is limited to a single reference case; additional case studies are needed to demonstrate model robustness.}

\textbf{Response:}
The authors were unable to complete validation for additional reference cases in the time provided, and this caveat is noted in the limitations section. If the reviewers believe this limitation is preventing publication of the paper, with additional time we can complete further validation.

\item \textbf{Reviewer 1, Comment 4.}

\emph{The manuscript lacks a comprehensive sensitivity analysis or uncertainty analysis of key design parameters.}

\textbf{Response:}
The design of experiments serves as a preliminary analysis of key design parameters. The authors have conducted a comprehensive sensitivity analysis but, due to the paper's length, have decided to publish it separately, so that the present manuscript can focus on the development and validation of the model rather than its results.

\item \textbf{Reviewer 1, Comment 5.}

\emph{Add a workflow diagram summarizing the complete MDOcean framework.}

\textbf{Response:}
The existing XDSM diagram shows the workflow. A new scope diagram has been added to contextualize this workflow within a broader WEC modeling context.

\item \textbf{Reviewer 1, Comment 6.}

\emph{Add a notation list.}

\textbf{Response:}
A glossary has been added.

\item \textbf{Reviewer 1, Comment 7.}

\emph{Several mathematical derivations are difficult to follow and require clearer explanations or simplification.}

\textbf{Response:}
% TODO: Add response.

\item \textbf{Reviewer 1, Comment 8.}

\emph{Eq.~1--38: I would not criticize the mathematics itself, but I would comment on clarity, notation consistency, assumptions, and reproducibility, as these are the more defensible reviewer observations.}

\textbf{Response:}
The assumption numbering and diagram intend to address this recommendation.

\item \textbf{Reviewer 1, Comment 9.}

\emph{Better highlight the multidisciplinary integration.}

\textbf{Response:}
% TODO: Add response.

\item \textbf{Reviewer 1, Comment 10.}

\emph{There should be references in the equations.}

\textbf{Response:}
Any equations without a reference were derived by the authors.

\item \textbf{Reviewer 1, Comment 11.}

\emph{State study objectives more explicitly at the end of the introduction.}

\textbf{Response:}
% TODO: Add response.

\item \textbf{Reviewer 1, Comment 12.}

\emph{Better highlight the open-source availability.}

\textbf{Response:}
% TODO: Add response.

\item \textbf{Reviewer 1, Comment 13.}

\emph{Discuss the limitations associated with the semi-analytical assumptions: deeper discussion on how they impact accuracy.}

\textbf{Response:}
% TODO: Add response.

\item \textbf{Reviewer 1, Comment 14.}

\emph{All key terms must be clearly defined at first use to avoid ambiguity and improve clarity.}

\textbf{Response:}
The glossary of acronyms and symbols addresses this recommendation.

\item \textbf{Reviewer 1, Comment 15.}

\emph{All parameters appearing in the equations need to be clearly defined and explained.}

\textbf{Response:}
The glossary of acronyms and symbols addresses this recommendation.

\item \textbf{Reviewer 1, Comment 16.}

\emph{The novelty of MDOcean compared to existing WEC simulation frameworks is not sufficiently highlighted.}

\textbf{Response:}
% TODO: Add response.

\item \textbf{Reviewer 1, Comment 17.}

\emph{The structural and economic models are relatively simplified and need stronger justification for their applicability.}

\textbf{Response:}
% TODO: Add response.

\item \textbf{Reviewer 1, Comment 18.}

\emph{The relevant papers should be reviewed and cited. [Several examples]}

\textbf{Response:}
Not addressed in the current revision.

\item \textbf{Reviewer 1, Comment 19.}

\emph{The limitations of the proposed framework and its applicability to other WEC configurations should be discussed in greater detail.}

\textbf{Response:}
% TODO: Add response.

\item \textbf{Reviewer 1, Comment 20.}

\emph{The abstract is too wordy. Word count is 254, limit is 250.}

\textbf{Response:}
% TODO: Add response.

\item \textbf{Reviewer 1, Comment 21.}

\emph{Put figures in their places.}

\textbf{Response:}
Figures are in their places in the source, and their rendering at the end of the document is a result of the provided Elsevier LaTeX template.

\item \textbf{Reviewer 1, Comment 22.}

\emph{What can I understand from Fig.~16 and 17? Too messy.}

\textbf{Response:}
% TODO: Add response.

% ----------------------------------------------------------------
% Reviewer 2
% ----------------------------------------------------------------

\item \textbf{Reviewer 2, Comment 1.}

\emph{In Section 3.3.4, a constant efficiency $\eta$ is used to compute average electrical power from average mechanical power. This appears to be a main limitation, as it is highlighted in a number of studies on loss-aware control of WECs that a constant efficiency may not be accurate for either power evaluation or control, given the oscillating operation of WECs. While an $\eta$ is selected to match a specific operating condition, the efficiency may change when the motion amplitude changes.}

\textbf{Response:}
The presented equations do include PTO dynamics (non-constant efficiency such as ohmic loss and inductive effects). The constant efficiency is used for validation against the RM3 report, since that is what the latter assumes.

\item \textbf{Reviewer 2, Comment 2.}

\emph{In Section 4, while you have presented validation results on the dynamic model versus WEC-Sim, please discuss the validation on the optimal control solution against an ideal benchmark such as MPC. Given that the proposed framework has a number of approximate treatments in constraint/nonlinearity handling and monochromatic simplification of irregular waves, the accuracy of each assumption, as well as how much they affect the final evaluation results, remains unclear.}

\textbf{Response:}
While we do not validate against MPC, prior results using the WecOptTool software have indicated that a reactive (PI) controller obtains 92% of the power of a high-order controller. The text describes this and it is a suitable level of fidelity for the intended early-stage design optimization.

\item \textbf{Reviewer 2, Comment 3.}

\emph{In Section 3.3.4, clearly define $\hat{e}$ and $\hat{q}$ in the main text. Are they complex frequency-domain coefficients? It would also be helpful to explicitly state here that the derivation focuses on a single frequency.}

\textbf{Response:}
% TODO: Add response.

\item \textbf{Reviewer 2, Comment 4.}

\emph{In Section 3.3.5, the second part is titled `constrained'' and the third part is `nonlinear controller.'' However, the third part also addresses a specific constraint (force limit). Clarify what specific constraints these two parts are intended to handle and use clearer terminology.}

\textbf{Response:}
% TODO: Add response.

\item \textbf{Reviewer 2, Comment 5.}

\emph{In Section 3.3.5, the `constrained'' part before Eq.~(18), `each metric is written as a quadratic product of the port variable vector with a matrix.'' Is this valid for all metrics that may need to be constrained? If a WEC has a velocity/displacement constraint, should the constrained metric be the port variable itself rather than the quadratic function? This goes back to the variable definition.}

\textbf{Response:}
% TODO: Add response.

\item \textbf{Reviewer 2, Comment 6.}

\emph{Throughout the derivation in the main text, a clear definition of describing function is missing, which makes it difficult to follow the content.}

\textbf{Response:}
A clear definition of describing function has been added in Section 3.3.1.

\item \textbf{Reviewer 2, Comment 7.}

\emph{In Section 3.3.5 and the corresponding appendix, I did not fully understand how the force saturation is applied. Are you simply scaling the force limit by $4/\pi$ and applying it on the force coefficient?}

\textbf{Response:}
% TODO: Add response.

\item \textbf{Reviewer 2, Comment 8.}

\emph{In Section 3.3.6, Eq.~24, are these coefficients fixed, or do they depend on the motion amplitude?}

\textbf{Response:}
They depend on the motion amplitude.

% ----------------------------------------------------------------
% Reviewer 3
% ----------------------------------------------------------------

\item \textbf{Reviewer 3, Comment 1.}

\emph{I believe the proposed framework will be useful from the practical viewpoint, but from the viewpoint of a research journal paper, I cannot find any innovative new findings or developments worth being published as a journal paper.}

\textbf{Response:}
The revised manuscript emphasizes the aspects which are novel. The non-novel aspects are still significant because they synthesize previously disparate models and comment on the coupling relations between modules.

\item \textbf{Reviewer 3, Comment 2.}

\emph{Conclusion should be shortened and should distinguish demonstrated results from intended future capabilities.}

\textbf{Response:}
The conclusion has been shortened.

\item \textbf{Reviewer 3, Comment 3.}

\emph{The manuscript requires substantial revision in terms of clarity, validation, novelty justification, and presentation before it can be considered for publication.}

\textbf{Response:}
The authors hope that the reorganization of the manuscript, particularly of Section 3.3, provides clarity.

\item \textbf{Reviewer 3, Comment 4.}

\emph{The benchmarking focuses mainly on computational speed, while accuracy comparisons should be expanded.}

\textbf{Response:}
Not addressed in the current revision.

\item \textbf{Reviewer 3, Comment 5.}

\emph{Several tables and figures are dense. Validation and benchmarking results would be easier to assess if the key errors and runtime comparisons were summarized in one compact table.}

\textbf{Response:}
Not addressed in the current revision.

\item \textbf{Reviewer 3, Comment 6.}

\emph{The manuscript should avoid relying too heavily on forthcoming companion papers for essential derivations, implementation details, or validation evidence.}

\textbf{Response:}
The authors have not included additional material in this paper due to competing concerns about the length of the manuscript. The companion papers are available as preprints for interested readers while their publication process is ongoing.

% ----------------------------------------------------------------
% Reviewer 4
% ----------------------------------------------------------------

\item \textbf{Reviewer 4, Comment 1.}

\emph{For the spar and coupling coefficient algebraic approximations, interpolation, and corrections for positive definiteness, explain how much they change the dynamic response or power estimates and how often corrections are needed.}

\textbf{Response:}
This can be seen by comparing the top and bottom rows of Figure~36.

\item \textbf{Reviewer 4, Comment 2.}

\emph{My main concern is the way the accuracy and validation are presented. Some results appear calibrated to RM3 reference data with scaling factors, while others are compared to WEC-Sim under matched modeling assumptions. These are both useful exercises, but they should not be presented in the same way as independent validation. The authors should more carefully distinguish validation, calibration, verification, and benchmarking, and should moderate claims about model accuracy and applicability.}

\textbf{Response:}
% TODO: Add response.

\item \textbf{Reviewer 4, Comment 3.}

\emph{Scale factors are reasonable for reproducing RM3, but raise an important question: how would these factors be selected for a new WEC design without reference data? Discuss this point and, if possible, show the main validation errors before and after scaling.}

\textbf{Response:}
% TODO: Add response.

\item \textbf{Reviewer 4, Comment 4.}

\emph{The structural module is semi-analytical, and FEA-based validation is still listed as future work. Avoid giving the impression that the structural module is already suitable for detailed structural design or certification. It is better presented as an early-stage screening tool.}

\textbf{Response:}
The intention is for the model to be used as an early-stage screening tool to explore large design spaces, rather than to perform late-stage performance optimization of an existing WEC.

\item \textbf{Reviewer 4, Comment 5.}

\emph{Use of equivalent regular waves and regular wave theory for storm cases is a major approximation. This can miss nonlinear wave effects and transient load peaks. The authors acknowledge this limitation, but the main text and conclusion should make it clearer that the storm load and survivability predictions are preliminary and should not replace higher-fidelity analysis.}

\textbf{Response:}
% TODO: Add response.

\item \textbf{Reviewer 4, Comment 6.}

\emph{Use of a JPD is useful for long-term performance estimates, but the dynamics are still based on equivalent regular wave approximations. Discuss how this affects annual average power, fatigue-related quantities, and controller performance, especially for broadbanded sea states.}

\textbf{Response:}
% TODO: Add response.

\item \textbf{Reviewer 4, Comment 7.}

\emph{Avoid implying that the framework is already applicable to all WEC archetypes---it is only demonstrated for RM3 two-body axisymmetric point absorbers. The applicability is tied to the assumptions of the hydrodynamic, structural, and economic modules.}

\textbf{Response:}
It is only the dynamics and control module that applies to all WEC archetypes.

\item \textbf{Reviewer 4, Comment 8.}

\emph{The abstract should mention the most important limitations more explicitly, especially the regular-wave approximation and the focus on a single WEC archetype.}

\textbf{Response:}
% TODO: Add response.

\item \textbf{Reviewer 4, Comment 9.}

\emph{Please use `validation,'' `calibration,'' `verification,'' and `benchmarking'' consistently throughout the manuscript.}

\textbf{Response:}
% TODO: Add response.

\item \textbf{Reviewer 4, Comment 10.}

\emph{The term ``near-numerical accuracy'' should be used carefully. The model is fast and useful, but some agreement is obtained after calibration.}

\textbf{Response:}
% TODO: Add response.

\item \textbf{Reviewer 4, Comment 11.}

\emph{The open-source information should be made easier to find, including repository link, version number, documentation status, and scripts for reproducing the main validation and benchmarking cases.}

\textbf{Response:}
Not addressed in the current revision.

\item \textbf{Reviewer 4, Comment 12.}

\emph{The economic model uses 2012 USD for consistency with RM3. This is acceptable, but it should be clearly stated whenever LCOE or cost values are discussed.}

\textbf{Response:}
Not addressed in the current revision.

\item \textbf{Reviewer 4, Comment 13.}

\emph{Agreement with RM3 should be discussed as calibration rather than independent validation. Clearly state which quantities are predicted without tuning and which quantities are matched after scaling factors.}

\textbf{Response:}
% TODO: Add response.

\item \textbf{Reviewer 4, Comment 14.}

\emph{Make clear what is included in the 151~ms runtime: single sea state or full JPD, one design evaluation or one module call, with/without hydrodynamic precomputation, data loading, and setup costs. This is important for readers who may want to compare MDOcean with BEM, WEC-Sim, or other optimization workflows.}

\textbf{Response:}
This is a full JPD of one design evaluation with all modules. It does not include setup costs. Hydrodynamic precomputation is not applicable because the hydrodynamics module is part of the simulation.

\end{enumerate}

\end{document}
